% ImpactKV M3 prefetch 汇报稿（中文）
% 编译：lualatex -interaction=nonstopmode PREFETCH_M3_REPORT_CN.tex
\documentclass[UTF8,a4paper,11pt,fontset=fandol]{ctexart}
\usepackage[margin=2.2cm]{geometry}
\usepackage{booktabs}
\usepackage{tabularx}
\usepackage{enumitem}
\usepackage[ruled,vlined,linesnumbered]{algorithm2e}
\usepackage{amsmath}
\usepackage{hyperref}
\usepackage{xcolor}
\usepackage{tikz}
\usepackage{graphicx}
\usepackage{caption}
\usepackage{listings}
\usetikzlibrary{arrows.meta,positioning,fit,calc,shapes.geometric}
\hypersetup{colorlinks=true,linkcolor=blue!50!black,urlcolor=blue!50!black}
\definecolor{ikvblue}{HTML}{2F6DB3}
\definecolor{ikvnavy}{HTML}{1B4F8A}
\definecolor{ikvteal}{HTML}{5BB3A7}
\definecolor{ikvorange}{HTML}{D96B2B}
\definecolor{ikvgray}{HTML}{6B6B6B}
\captionsetup{font=small,labelfont=bf}
\setlist[itemize]{leftmargin=1.4em,itemsep=0.25em,topsep=0.3em}
\setlist[enumerate]{leftmargin=1.6em,itemsep=0.25em,topsep=0.3em}
\title{\textbf{ImpactKV 模板引导 Prefetch（M3）算法汇报}}
\author{ImpactKV / ASPLOS 2027}
\date{2026-08-26}
\begin{document}
\maketitle
\begin{quote}\small
本报告说明 \textbf{M3 模板引导 prefetch}：它做什么、不做什么、如何与 M0/M2 接口、
顺序 one-token 评测上为什么不是加速、以及 headline 为何关掉它。
数字来自冻结战役，禁止改作业 137185 / 96092 的 \texttt{RESULT.json}。
\end{quote}

\tableofcontents
\vspace{0.6em}
\noindent\rule{\linewidth}{0.4pt}

\section{一句话结论}

M3 回答的是 \textbf{何时、把已经 admit 的页放到哪一层}（device / host），
\textbf{不回答哪些 span 可以拷}。Admit 是 M0，拷贝是 M2。
顺序 1-token、source 已在 device 时，没有重叠窗口，prefetch 不能比「已经持有的拷贝」更快。
Headline 7B（作业 137185，cache-ready $1.492\times$）把 M1 和 M3 \textbf{关掉}，
主表不能被 radix 或 prefetch 领功。

\begin{center}
\small
\begin{tabular}{llp{0.42\textwidth}}
\toprule
\textbf{模块} & \textbf{Headline 137185} & \textbf{职责} \\
\midrule
M0 编译器 & 开 & 离线 PLAN：单文件、版本有效、token-ID 相同、$\Delta\neq 0$ \\
M1 前缀 & \textbf{关} & radix LCP（$\Delta=0$），不拷岛 \\
M2 拷贝 & 开 & $K\leftarrow R_\Delta K$，$V$ 原样；失败则整岛 Dense \\
M3 prefetch & \textbf{关} & 对已 admit 的页发驻留 hint；miss $\to$ 仍持有的 M2 \\
\bottomrule
\end{tabular}
\end{center}

\section{问题：驻留 $\neq$ 准入}

编码 agent 的 planner / implementer / tester / reviewer 会在滚动历史里
反复读同一仓库文件。文件岛的 token ID 相同，位置不同（$\Delta\neq 0$）。
这带来两个独立问题：

\begin{enumerate}
  \item \textbf{Admit（M0）}：哪些 span 允许 true-lossy 拷贝。
  \item \textbf{Residency（M3）}：这些页现在在 host 还是 device；下一次用之前要不要提前搬。
\end{enumerate}

把两件事绑在一个开关里会出两类错：
(i)~prefetch 把不该拷的 tool log 搬上 GPU；
(ii)~prefetch miss 把本来已经 admit 的岛改成 Dense，等于关掉 M2。
第二条在 7 组 dual-island 战役上真实发生过：combined 相对 dual 变成 $0.528\times$（7 个 hint、4 次 miss 后 Dense）。
\textbf{正确语义：miss 必须落到仍持有的 M2 拷贝，不得 Dense 重算合法岛。}

\section{接口：M3 不扩大拷贝集合}

页键与 M2 相同：
\[
(\textit{source-prefix hash},\; \textit{content hash},\; \Delta).
\]
Hint 只携带已存在的键、目标层、相对截止时间和优先级。\textbf{不旋转 $K$}，
\textbf{不改 admit 集合}，\textbf{不估 Attention}。

实现上分两层编译器（\texttt{kvcomm\_prefetch/template\_hints.py}）：

\begin{itemize}
  \item \texttt{compile\_template\_prefetch\_hints}：优先级 $=$ 声明的剩余 use 次数（oracle，评测里不用它当 online 信号）。
  \item \texttt{compile\_next\_island\_prefetch\_hints}：\textbf{线上实际用的}。优先级 $=$ 编码协议里的 \textbf{later-roles}（planner 读完通常还有 3 个下游角色），\textbf{不是} 未来 \texttt{target\_uses}。
\end{itemize}

Skip 条件（任一成立则不发 hint）：

\begin{itemize}
  \item 未过 M0 门（非单文件、token-ID 不匹配、版本失效、$\Delta=0$ 的 MIDDLE）。
  \item \texttt{later\_roles\_in\_protocol} $=0$（协议上没有下游再读）。
  \item 页已经在 DEVICE \textbf{且} 下一个 use 就是当前顺序请求：\texttt{no\_overlap\_window}。
\end{itemize}

最后一条解释了作业 119795 上 \textbf{Hints $=0$}：顺序 1-token、岛已在 device，没有「提前搬」的时间窗。

\section{算法}

\subsection{Hint 编译（next-island / later-roles）}

\begin{algorithm}[ht]
\SetAlFnt{\footnotesize}
\DontPrintSemicolon
\KwIn{已观察的 PREFIX / MIDDLE 岛（协议信号，不含未来 target\_uses）}
\KwOut{按优先级排序的 \texttt{KVPrefetchHint}}
\BlankLine
\For{each observation $o$}{
    \If{$o$ 未过 admit 门}{\textbf{continue}}
    \If{$o$ 已在 DEVICE 且下一 use 即当前请求}{\textbf{continue}\tcp*{无重叠窗}}
    \If{$o$ 是 MIDDLE 且 later-roles $\le 0$}{\textbf{continue}}
    发 hint：$(\mathrm{key},\;\mathrm{DEVICE},\;\mathrm{deadline},\;\mathrm{priority}{=}\text{later-roles})$\;
}
按 $(-\mathrm{priority},\;\mathrm{key})$ 排序去重\;
\caption{编译 next-island prefetch hints（不扩大 admit 集）}
\end{algorithm}

\subsection{调度与加载}

\texttt{AsyncKVPrefetchScheduler} 是截止时间 / 优先级堆（相对 \texttt{deadline\_s} 从提交时刻起算）。
工作线程到期则标记 \texttt{expired}，否则交给 \texttt{KVPrefetchCoordinator}：

\begin{enumerate}
  \item 关核或关 prefetch $\to$ 直接 disabled。
  \item 按 $(\mathrm{deadline},\;-\mathrm{priority},\;\mathrm{key})$ 去重排序。
  \item \texttt{store.lookup(key)}：没有页 $\to$ \texttt{store\_miss}（\textbf{不创造}新岛）。
  \item 已在目标层 $\to$ pin；否则 \texttt{ensure\_resident} 再 pin。
\end{enumerate}

\texttt{MiddleKVPrefetchAPI.export\_middle\_kv} 可以把 device 切片备份到 host，
供稍后 prefetch 回 device。这是驻留搬运，不是第二次 admit。

\subsection{Miss 语义（必须）}

\begin{center}
\texttt{resolve\_copy\_source\_after\_prefetch(owned, prefetched)}
$=\;$ \texttt{prefetched} 若非空，否则 \texttt{owned}。
\end{center}

等待 hint 超时或 \texttt{MiddleKVPrefetchError} 时记 \texttt{template\_prefetch\_miss}，
然后仍用 owned handle 走 M2。返回空才允许 Dense。
代码注释写明：7 组 combined 的 4/28 fallback 把 dual 的 $8.49\times$ 打成 $4.48\times$，就是 miss 误走 Dense。

\section{运行时路径}

\begin{center}
\begin{tikzpicture}[
  font=\sffamily\scriptsize,
  >=Stealth,
  box/.style={draw=#1, rounded corners=2pt, align=center,
    inner sep=3.4pt, minimum height=0.78cm, fill=white},
]
\node[box=ikvteal, text width=2.15cm] (m0) at (0,1.35) {M0 PLAN\\已 admit 的岛};
\node[box=ikvblue, text width=2.35cm] (src) at (-3.2,0) {Source Dense\\写入租约池};
\node[box=ikvorange, text width=2.55cm] (hint) at (0,0) {M3 next-island\\later-roles hint};
\node[box=ikvgray, text width=2.15cm] (sched) at (3.2,0) {调度器\\deadline / 优先级};
\node[box=ikvnavy, text width=2.55cm] (tgt) at (0,-1.45) {Target：Dense 前缀 $\to$ M2 拷贝 $\to$ Dense 剩余};
\draw[->, ikvgray, thick] (m0) -- (src);
\draw[->, ikvgray, thick] (m0) -- (hint);
\draw[->, ikvorange, thick] (hint) -- (sched);
\draw[->, ikvblue, thick] (src) -- (tgt);
\draw[->, ikvorange] (sched.south) -- node[font=\tiny, text=ikvorange, right] {hit: 已在 device} (tgt.east);
\node[font=\sffamily\tiny, text=ikvgray, align=center] at (3.2,-1.45)
  {miss: 仍用 owned M2\\不得 Dense 合法岛};
\end{tikzpicture}
\end{center}
\captionof{figure}{M3 不改 PLAN。它只给已存在的页发驻留提示。Headline 战役此路径关闭。}

环境开关（\texttt{template\_prefetch\_modes.py}）：

\begin{center}
\small
\begin{tabular}{lccc}
\toprule
\textbf{模式} & \texttt{KVCOMM} & \texttt{LOSSY} & \texttt{PREFETCH} \\
\midrule
dense & 0 & 0 & 0 \\
prefix\_only / lossy\_only / dual & 1 & 1 & 0 \\
combined & 1 & 1 & 1 \\
\bottomrule
\end{tabular}
\end{center}

\texttt{combined} 才打开 host overlap 与 \texttt{prefetch\_spill\_device}。
Manifest 里 \texttt{prefetch\_deadline\_s}$=30$、\texttt{prefetch\_wait\_s}$=90$。
\texttt{prefix\_only} 把 \texttt{copy\_middle=False}，只留 radix 前缀。

\section{实验（冻结数字，禁止手改 RESULT）}

\subsection{作业 119795：235 组顺序 1-token（30B PLAN，附录）}

同一 235 组、$705$ 对。Hints 是 next-island / later-roles，不是 remaining uses。
顺序下下一岛已在 DEVICE，hint 被 \texttt{no\_overlap\_window} 跳过。

\begin{center}
\small
\begin{tabular}{lrrrrr}
\toprule
\textbf{臂} & \textbf{相对本战役 Dense} & \textbf{拷贝} & \textbf{Spill} & \textbf{Hints} & \textbf{Miss} \\
\midrule
Dense & $1.000\times$ & --- & --- & --- & --- \\
coding（仅拷贝） & $1.390\times$ & $1684/1684$ & $0$ & $0$ & $0$ \\
prefetch-only & $0.996\times$ & $0$ & $0$ & $0$ & $0$ \\
combined & $1.392\times$ & $1684/1684$ & $0$ & $0$ & $0$ \\
\bottomrule
\end{tabular}
\end{center}
\captionof{table}{作业 119795。combined $\approx$ coding，是开销不是 prefetch 加速。
prefetch-only 不是拷贝赢。\textbf{不要}把此表当成 7B 主表 $1.492\times$。}

要点：

\begin{itemize}
  \item Dual 前缀$+$拷贝相对 Dense \textbf{不是} prefetch 数字。
  \item 119795 \textbf{不是} 96092（96092 是 30B 拷贝附录，$1.375\times$，prefetch 关）。
  \item 论文 \texttt{tab:template-prefetch} 停在 119795。
\end{itemize}

\subsection{7 组 dual-island（作业 124825--124829）}

7 组、21 对。prefix-only $1.155\times$，lossy-only $4.526\times$，dual $8.490\times$。
\texttt{combined} 发了 7 个 hint、4 次 miss，随后 \textbf{Dense 了合法岛}（24/28 拷、4 fallback），
相对 dual 为 $0.528\times$。这是模块 bug，不是「prefetch 更慢」的算法结论。
当前 \texttt{kvcomm\_exact.py} 已改成 miss $\to$ owned M2。验证需要\textbf{新目录}重跑，
\textbf{禁止覆盖}上述冻结 RESULT。

\section{为何 Headline 关掉 M3}

顺序 cache-ready 协议（source KV 已在、解 1 token、对同一引擎 Dense）下：

\begin{itemize}
  \item 下一岛已经在 device，prefetch 没有重叠窗。
  \item 打开 M3 却把增益算给 coding-aware copy，审稿人会杀主表。
  \item 7 组上 miss$\to$Dense 会把拷贝核弄脏。
\end{itemize}

因此作业 137185：prefetch off、prefix off、1684/1684、0 fallback、$1.492\times$。
M3 的论文位置是\textbf{附录驻留检查}，不是速度声称。

\section{代码地图}

\begin{center}
\small
\begin{tabular}{>{\ttfamily}p{0.58\textwidth}p{0.36\textwidth}}
\toprule
\multicolumn{1}{l}{\textbf{路径}} & \textbf{角色} \\
\midrule
python/.../kvcomm\_prefetch/template\_hints.py & later-roles / next-island 编译 \\
python/.../kvcomm\_prefetch/scheduler.py & 截止时间堆、过期 \\
python/.../kvcomm\_prefetch/coordinator.py & lookup / load / pin \\
python/.../kvcomm\_prefetch/middle\_kv.py & host 导出与回装 \\
python/.../kvcomm\_exact.py & miss $\to$ owned M2 \\
benchmark/.../template\_prefetch\_modes.py & 五模式开关 \\
benchmark/.../run\_swebench\_template\_prefetch.py & 战役入口 \\
\bottomrule
\end{tabular}
\end{center}

单测（无 GPU）：
\texttt{kvcomm\_prefetch/test\_*.py}，含
\texttt{test\_prefetch\_store\_miss\_does\_not\_copy}、
\texttt{test\_resolve\_copy\_source\_after\_prefetch\_keeps\_owned\_on\_miss}。

\section{汇报时建议强调的三句话}

\begin{enumerate}
  \item Prefetch 是\textbf{驻留}，不是\textbf{准入}；hint 的键必须已经在 M0 白名单里。
  \item 信号来自编码协议的 later-roles / next-island，\textbf{不是} remaining $N$-use，也不是 Attention。
  \item 顺序 1-token 上 combined $\approx$ copy（$1.392\times$ vs $1.390\times$），不是加速；
        headline 关掉 M3，避免主表被驻留模块领功。
\end{enumerate}

\section{不要说的}

\begin{itemize}
  \item 不要说 M3 贡献了 $1.492\times$。
  \item 不要把 119795 和 96092、137185 写成一个 official method。
  \item 不要报 N=4，不要写 SOTA，不要把 one-token agreement 叫 Accuracy。
  \item 不要把 7 组 $0.528\times$ 解释成「prefetch 算法更慢」而不提 miss$\to$Dense bug。
\end{itemize}

\vspace{0.8em}
\noindent\small
冻结战役：119795（\texttt{impactkv\_swebench\_template\_prefetch\_nextisland\_20260821}）、
7 组 dual-island（\texttt{...\_7b\_dualisland\_20260822}）。
引擎分支 \texttt{integration/template-prefetch-swebench}。
接手说明见仓库根目录 \texttt{HANDOFF.md}。

\end{document}
